\documentclass[A4,svgnames,9pt,aspectratio=169]{beamer}
%% document options:
%% - aspectratio = { 43, 169, 1610 }
%% - utf8
%%

%%
%% insert list of packages
%%

%% \usepackage[english,french]{babel}

\hypersetup{
   allcolors=rouge_inria,
   pdfauthor   = {firstname lastname},
   pdftitle    = {\@title},
   pdfsubject  = {Resum\'{e} of everything},
   pdfkeywords = {firstname~lastname, curriculum vit\ae{}}
}

%%%%%%%%%%%%%%%%%%%%%%%%%%%%%%%%%%%%%%%%%%%%%%%%%%%%%%%
%%
%%%%%%%%%%%%%%%%%%%%%%%%%%%%%%%%%%%%%%%%%%%%%%%%%%%%%%%

\title[titrecourt]{Amélioration des performances I/O de Gysela}
\subtitle{Stage de M1, 2026}
\date[00/07/2026]{date long}
\author[A. et al.]{Alexandre Lou-Poueyou}

\usetheme{inria}

\begin{document}

%%%%%%%%%%%%%%%%%%%%%%%%%%%%%%%%%%%%%%%%%%%%%%%%%%%%%%%
%%
%%%%%%%%%%%%%%%%%%%%%%%%%%%%%%%%%%%%%%%%%%%%%%%%%%%%%%%

\frame{\titlepage}

%%%%%%%%%%%%%%%%%%%%%%%%%%%%%%%%%%%%%%%%%%%%%%%%%%%%%%%

% Le titre des planches de sommaire est \contentsname, sa valeur
% est fixée ici à "Sommaire" par défaut.
\renewcommand{\contentsname}{Sommaire}

% Si le package babel est utilisé la valeur
% prend celle par défaut de la langue sélectionnée.
% Pour imposer une valeur, y compris lors de l'utilisation de
% babel, faire attention à redéfinir cette variable _après_
% un \selectlanguage{}

% \selectlanguage{french}
% \selectlanguage{english}

\frame{\tocpage}

%%%%%%%%%%%%%%%%%%%%%%%%%%%%%%%%%%%%%%%%%%%%%%%%%%%%%%%
 
\section{Introduction}
\frame{\sectionpage}

%%%%%%%%%%%%%%%%%%%%%%%%%%%%%%%%%%%%%%%%%%%%%%%%%%%%%%%

\begin{frame}{Présentation Gysela}
  \begin{block}{Gysela}
    \begin{itemize}
      \item Code de simulation de turbulence dans les plasmas de fusion.
      \item Modèle de Vlasov gyro-kinétique 5D.
      \item C++/Fortran.
    \end{itemize}
  \end{block}

  \begin{block}{Problématique}
    \begin{itemize}
      \item Checkpoint régulier nécessaire pour éviter la perte de données en cas de crash.
      \item Bloque la simulation pendant l'écriture.
      \item Objectif : réduire le temps d'écriture.
    \end{itemize}
  \end{block}
\end{frame}

\subsection{Problématique I/O}

\begin{frame}{Problématique I/O}
  \begin{block}{Checkpointing}
    \begin{itemize}
      \item Grosse simulation nécessitant beaucoup de nœuds de calcul.
      \item Checkpoint régulier pour éviter la perte de données.
      \item Bloque la simulation pendant l'écriture.
    \end{itemize}
  \end{block}
\end{frame}

%%%%%%%%%%%%%%%%%%%%%%%%%%%%%%%%%%%%%%%%%%%%%%%%%%%%%%%

\section{Objectifs}
\frame{\sectionpage}

%%%%%%%%%%%%%%%%%%%%%%%%%%%%%%%%%%%%%%%%%%%%%%%%%%%%%%%

\begin{frame}{Objectifs du stage}
  \begin{block}{Objectifs}
    \begin{itemize}
      \item Analyser les performances I/O actuelles de Gysela.
      \item Implémenter des techniques d'écriture asynchrone.
      \item Évaluer les gains de performance obtenus.
    \end{itemize}
  \end{block}
\end{frame}

%%%%%%%%%%%%%%%%%%%%%%%%%%%%%%%%%%%%%%%%%%%%%%%%%%%%%%%

\section{Environnement d'expérimentation}
\frame{\sectionpage}

%%%%%%%%%%%%%%%%%%%%%%%%%%%%%%%%%%%%%%%%%%%%%%%%%%%%%%%

\begin{frame}{Environnement d'expérimentation}
    \begin{block}{}
      \begin{itemize}
        \item Calcul effectué sur Plafrim : Noeuds bora (44 nœuds).
        \begin{itemize}
          \item CPU : 2x 18-core Cascade Lake Intel Xeon Skylake Gold 6240 @ 2.6 GHz.
          \item Mémoire : 192 GB (5.3 GB/core) @ 2993 MT/s.
          \item Network : Omnipath 100Gbit/s \& Ethernet 10Gbit/s.
        \end{itemize}
        \item Dimension de la simulation : 128x128x128x64x64 (5D).
        \item Taille du checkpoint : 540Gb.
        \item Système de fichiers : BeeGFS (avec 8 OST maximum) et un chunksize de 512KB.
      \end{itemize}
    \end{block}
\end{frame}

\begin{frame}{Outils de mesure des performances}
    \begin{block}{IOPS}
        \begin{itemize}
            \item Indication des paramètres à faire varier
            \item Extraction de métriques
            \begin{itemize}
              \item Temps d'écriture (metrics.time\_write)
              \item Temps total (metrics.total\_time)
              \item Utilisation de la bande passante (metrics.bandwidth)
            \end{itemize}
        \end{itemize}
    \end{block}
\end{frame}

%%%%%%%%%%%%%%%%%%%%%%%%%%%%%%%%%%%%%%%%%%%%%%%%%%%%%%% 

\section{Conclusion}
\begin{frame}{Conclusion}

    {\tiny         The quick brown fox jumps over the lazy dog}\\
    {\scriptsize   The quick brown fox jumps over the lazy dog}\\
    {\footnotesize The quick brown fox jumps over the lazy dog}\\
    {\small        The quick brown fox jumps over the lazy dog}\\
    {\normalsize   The quick brown fox jumps over the lazy dog}\\
    {\large        The quick brown fox jumps over the lazy dog}\\
    {\Large        The quick brown fox jumps over the lazy dog}\\
    {\LARGE        The quick brown fox jumps over the lazy dog}\\
    {\huge         The quick brown fox jumps over the lazy dog}\\
    {\Huge         The quick brown fox jumps over the lazy dog}

\end{frame}

%%%%%%%%%%%%%%%%%%%%%%%%%%%%%%%%%%%%%%%%%%%%%%%%%%%%%%%

%% Le texte est modifiable en changeant \thankyou
%% \renewcommand{\thankyou}{Thank You.}
\frame{\merci}

%%%%%%%%%%%%%%%%%%%%%%%%%%%%%%%%%%%%%%%%%%%%%%%%%%%%%%% 

\end{document}

%%%%%%%%%%%%%%%%%%%%%%%%%%%%%%%%%%%%%%%%%%%%%%%%%%%%%%%
%%
%%%%%%%%%%%%%%%%%%%%%%%%%%%%%%%%%%%%%%%%%%%%%%%%%%%%%%%

